\section{The Boltzmann target is an exact Gaussian mixture}\label{app:gmm}

The sampler state $\boldsymbol{\xi}$ is a free point in $\mathbb{R}^{d}$, where $d$
is the number of PCA coordinates. Each of the $K$ stored memories
$\mathbf{m}_1,\ldots,\mathbf{m}_K$ has length one and carries a multiplicity weight
$r_k > 0$. The state itself has no fixed length, because nothing in the update
rescales it back onto the unit sphere. It can land anywhere, which is why the answer
below is a distribution over all of $\mathbb{R}^{d}$.
Where does the sampler settle if we run it for a long time? For the underlying
continuous-time dynamics the answer is the Boltzmann distribution
$p_\beta(\boldsymbol{\xi})\propto e^{-\beta E_r(\boldsymbol{\xi})}$, built from the
multiplicity-weighted Hopfield energy
\begin{equation}
  E_r(\boldsymbol{\xi})
  = \tfrac12\lVert\boldsymbol{\xi}\rVert^2
    - \tfrac1\beta\log\sum_{k=1}^{K} r_k\, e^{\beta\, \mathbf{m}_k^\top\boldsymbol{\xi}},
  \qquad \lVert\mathbf{m}_k\rVert = 1 .
\end{equation}
This section shows that distribution is exactly a mixture of Gaussians. The update
we actually run, Eq.~\MainEqRef{eq:weighted-ula}, takes finite steps, so it only
comes close to this distribution rather than landing on it
exactly~\cite{durmusMoulinesULA2017}.

The proof takes three steps. First we exponentiate, which clears the logarithm and
leaves one sum over the memories. Then we complete the square, which puts each term
in the shape of a Gaussian. Then we normalize, which fixes the constants and gives
the mixture weights.

It helps to write down the shape we are aiming for first. A Gaussian in $d$
dimensions with mean $\boldsymbol{\mu}$ and the same variance $\sigma^{2}$ in every
direction has density
\begin{align}
  \mathcal{N}(\boldsymbol{\xi};\,\boldsymbol{\mu},\,\sigma^{2}\mathbf{I})
  &= \left(2\pi\sigma^{2}\right)^{-d/2}
     \exp\!\left[-\frac{1}{2\sigma^{2}}
       \lVert\boldsymbol{\xi}-\boldsymbol{\mu}\rVert^{2}\right]
     \notag\\[4pt]
  \mathcal{N}(\boldsymbol{\xi};\,\boldsymbol{\mu},\,\beta^{-1}\mathbf{I})
  &= \left(\tfrac{\beta}{2\pi}\right)^{\!d/2}
     \exp\!\left[-\tfrac\beta2
       \lVert\boldsymbol{\xi}-\boldsymbol{\mu}\rVert^{2}\right] .
  \label{eq:si-normal}
\end{align}
The second line is the first with $\sigma^{2}$ set to $\beta^{-1}$, the value we
need below. The only place $\boldsymbol{\xi}$ appears is inside
$\lVert\boldsymbol{\xi}-\boldsymbol{\mu}\rVert^{2}$, its squared distance to the
mean. So whenever we see something of the form
$\exp[-\tfrac\beta2\lVert\boldsymbol{\xi}-\mathbf{c}\rVert^{2}]$, we are looking at
a Gaussian centered at $\mathbf{c}$ with variance $\beta^{-1}$, apart from the
constant $(\beta/2\pi)^{d/2}$ in front of it. Everything below is getting each term
into that form with $\mathbf{c} = \mathbf{m}_k$. That is what makes each memory the
center of a Gaussian, and $\beta$ the thing that sets its width.

\paragraph{Exponentiation.}
Multiply the energy by $-\beta$. That $-\beta$ cancels the $-\tfrac1\beta$ sitting
in front of the logarithm, so the log term is left with nothing multiplying it:
\begin{align}
  -\beta E_r(\boldsymbol{\xi})
  &= -\beta\left[\tfrac12\lVert\boldsymbol{\xi}\rVert^2
     - \tfrac1\beta\log\sum_{k=1}^{K} r_k\,
       e^{\beta\, \mathbf{m}_k^\top\boldsymbol{\xi}}\right]
     \notag\\[2pt]
  &= -\tfrac\beta2\lVert\boldsymbol{\xi}\rVert^2
     + \log\sum_{k=1}^{K} r_k\, e^{\beta\, \mathbf{m}_k^\top\boldsymbol{\xi}} .
\end{align}
Now raise $e$ to both sides. A sum in the exponent becomes a product of two
factors, and $e^{\log S} = S$ cancels the logarithm in the second one, which frees
the sum:
\begin{align}
  e^{-\beta E_r(\boldsymbol{\xi})}
  &= e^{-\frac\beta2\lVert\boldsymbol{\xi}\rVert^2}\;
     \exp\!\left[\log\sum_{k=1}^{K} r_k\,
       e^{\beta\, \mathbf{m}_k^\top\boldsymbol{\xi}}\right]
     \notag\\[2pt]
  &= e^{-\frac\beta2\lVert\boldsymbol{\xi}\rVert^2}
     \sum_{k=1}^{K} r_k\, e^{\beta\, \mathbf{m}_k^\top\boldsymbol{\xi}}
     \notag\\[2pt]
  &= \sum_{k=1}^{K} r_k\,
     \exp\!\left[-\tfrac\beta2\lVert\boldsymbol{\xi}\rVert^2
       + \beta\, \mathbf{m}_k^\top\boldsymbol{\xi}\right] .
\end{align}
The last line pulls the $e^{-\frac\beta2\lVert\boldsymbol{\xi}\rVert^{2}}$ factor
inside the sum, which is fine because it does not depend on $k$. From here we work
on one term at a time, and the job is to make its exponent look like the one in
Eq.~\ref{eq:si-normal}.

\paragraph{Completing the square.}
We want the exponent
$-\tfrac\beta2\lVert\boldsymbol{\xi}\rVert^{2}
 + \beta\,\mathbf{m}_k^\top\boldsymbol{\xi}$
to look like $-\tfrac\beta2\lVert\boldsymbol{\xi}-\mathbf{m}_k\rVert^{2}$. It almost
does. Both start with $-\tfrac\beta2\lVert\boldsymbol{\xi}\rVert^{2}$, but in ours
the memory sits outside a square instead of inside one. The quickest way to see the
difference is to start from what we want and expand it:
\begin{align}
  -\tfrac\beta2\lVert\boldsymbol{\xi}-\mathbf{m}_k\rVert^2
  &= -\tfrac\beta2\left[\lVert\boldsymbol{\xi}\rVert^2
     - 2\,\mathbf{m}_k^\top\boldsymbol{\xi}
     + \lVert\mathbf{m}_k\rVert^2\right]
     \notag\\[2pt]
  &= \underbrace{-\tfrac\beta2\lVert\boldsymbol{\xi}\rVert^2
     + \beta\,\mathbf{m}_k^\top\boldsymbol{\xi}}_{\text{exponent of the }k\text{th term}}
     \;-\; \tfrac\beta2\lVert\mathbf{m}_k\rVert^2 .
\end{align}
The underbraced part is our exponent. So our exponent equals
$-\tfrac\beta2\lVert\boldsymbol{\xi}-\mathbf{m}_k\rVert^{2}
 + \tfrac\beta2\lVert\mathbf{m}_k\rVert^{2}$, which is the shape we want plus a
constant. Putting that into every term, then splitting the constant off as its own
factor:
\begin{align}
  e^{-\beta E_r(\boldsymbol{\xi})}
  &= \sum_{k=1}^{K} r_k\,
     \exp\!\left[-\tfrac\beta2\lVert\boldsymbol{\xi}-\mathbf{m}_k\rVert^2
       + \tfrac\beta2\lVert\mathbf{m}_k\rVert^2\right]
     \notag\\[2pt]
  &= \sum_{k=1}^{K} r_k\, e^{\frac\beta2\lVert\mathbf{m}_k\rVert^2}\,
     e^{-\frac\beta2\lVert\boldsymbol{\xi}-\mathbf{m}_k\rVert^2} .
\end{align}
The last factor is now the Gaussian shape from Eq.~\ref{eq:si-normal}, with
$\mathbf{c} = \mathbf{m}_k$. So term $k$ is a Gaussian centered on memory $k$ with
variance $\beta^{-1}$: a large $\beta$ makes it a tight bump around that memory, a
small $\beta$ a broad one. All that is left is the constants.

\paragraph{Unit norm and normalization.}
This is the step where unit norm matters. Because $\lVert\mathbf{m}_k\rVert = 1$
for every $k$, the leftover constant $e^{\frac\beta2\lVert\mathbf{m}_k\rVert^2}$ is
the same number $e^{\beta/2}$ in every term, so it comes out of the sum. Each
remaining
factor still needs the constant $(\beta/2\pi)^{d/2}$ to be a full Gaussian density,
so we put it in and divide it back out, which is the $(2\pi/\beta)^{d/2}$ in front:
\begin{align}
  e^{-\beta E_r(\boldsymbol{\xi})}
  &= e^{\beta/2}\sum_{k=1}^{K} r_k\,
     e^{-\frac\beta2\lVert\boldsymbol{\xi}-\mathbf{m}_k\rVert^2}
     \notag\\[2pt]
  &= e^{\beta/2}\left(\tfrac{2\pi}{\beta}\right)^{\!d/2}
     \sum_{k=1}^{K} r_k\,
     \mathcal{N}(\boldsymbol{\xi};\,\mathbf{m}_k,\,\beta^{-1}\mathbf{I}) .
\end{align}
To turn the $\propto$ into an $=$, we need the total $Z$, the integral over all of
$\mathbb{R}^d$. Every Gaussian density integrates to one, so the integral leaves
nothing behind but the sum of the $r_j$:
\begin{align}
  Z &= \int_{\mathbb{R}^d} e^{-\beta E_r(\boldsymbol{\xi})}\,
       \mathrm{d}\boldsymbol{\xi}
     \notag\\[2pt]
    &= e^{\beta/2}\left(\tfrac{2\pi}{\beta}\right)^{\!d/2}
       \sum_{j=1}^{K} r_j
       \underbrace{\int_{\mathbb{R}^d}
         \mathcal{N}(\boldsymbol{\xi};\,\mathbf{m}_j,\,\beta^{-1}\mathbf{I})\,
         \mathrm{d}\boldsymbol{\xi}}_{=\;1}
     \notag\\[2pt]
    &= e^{\beta/2}\left(\tfrac{2\pi}{\beta}\right)^{\!d/2}
       \sum_{j=1}^{K} r_j .
\end{align}
Dividing one by the other, $e^{\beta/2}$ and $(2\pi/\beta)^{d/2}$ sit on both the
top and the bottom, so they cancel and never matter:
\begin{align}
  p_\beta(\boldsymbol{\xi})
  &= \frac{e^{-\beta E_r(\boldsymbol{\xi})}}{Z}
   = \frac{e^{\beta/2}\left(\tfrac{2\pi}{\beta}\right)^{\!d/2}
           \sum_{k=1}^{K} r_k\,
           \mathcal{N}(\boldsymbol{\xi};\,\mathbf{m}_k,\,\beta^{-1}\mathbf{I})}
          {e^{\beta/2}\left(\tfrac{2\pi}{\beta}\right)^{\!d/2}
           \sum_{j=1}^{K} r_j}
     \notag\\[4pt]
  &= \sum_{k=1}^{K}
     \frac{r_k}{\sum_{j=1}^{K} r_j}\,
     \mathcal{N}(\boldsymbol{\xi};\,\mathbf{m}_k,\,\beta^{-1}\mathbf{I}) .
\end{align}
Writing $w_k = r_k / \sum_j r_j$ for the normalized multiplicities, this is
\begin{equation}
  \boxed{\;
  p_\beta(\boldsymbol{\xi})
  = \sum_{k=1}^{K} w_k\, \mathcal{N}(\boldsymbol{\xi};\,\mathbf{m}_k,\,\beta^{-1}\mathbf{I}),
  \qquad w_k = \frac{r_k}{\sum_j r_j} \; }
\end{equation}
an exact mixture of Gaussians, one per stored memory, each centered on its own
memory with variance $\beta^{-1}$, and weighted by its multiplicity divided by the
total. If the memories had not been unit norm, the
$e^{\frac\beta2\lVert\mathbf{m}_k\rVert^2}$ would differ from term to term and would
not come out of the sum, so the weights would be
$r_k\,e^{\frac\beta2\lVert\mathbf{m}_k\rVert^2}$ instead. The result would still be
a Gaussian mixture, just not with these weights.

\paragraph{Corollaries.}
\emph{(i) Exact sampling.} Draw $k\sim\mathrm{Categorical}(w)$ and
$\boldsymbol{\xi}=\mathbf{m}_k+\beta^{-1/2}\mathbf{z}$, $\mathbf{z}\sim\mathcal{N}(\mathbf{0},\mathbf{I})$;
no Markov chain is required. The continuous-time Langevin diffusion has
$p_\beta$ as its stationary law, while finite-step ULA is subject to
discretization and finite-time error.
\emph{(ii) Exact designated control.} Let the designated subset $B$ hold
$K_{\mathrm{des}}$ of the $K$ patterns, with the remaining
$K_{\mathrm{bg}} = K - K_{\mathrm{des}}$ forming the background. The probability
mass on $B$ is $\sum_{k\in B} w_k = \sum_{k\in B} r_k / \sum_j r_j
\equiv f_{\mathrm{eff}}$. Setting $r_k=\rho$ on $B$ and $1$ otherwise gives
$f_{\mathrm{eff}} = \rho K_{\mathrm{des}} / (\rho K_{\mathrm{des}} +
K_{\mathrm{bg}})$, with no equal-similarity assumption.
\emph{(iii) Score.} $\nabla_{\boldsymbol{\xi}}\log p_\beta
= -\beta\nabla E_r
= \beta\bigl[\mathbf{X}\,\mathrm{softmax}(\beta\mathbf{X}^\top\boldsymbol{\xi}+\log\mathbf{r})
  - \boldsymbol{\xi}\bigr]$, whose Euler-Maruyama discretization is
Eq.~\MainEqRef{eq:weighted-ula}.

\paragraph{Exact-equilibrium benchmark.}
We compared the finite-step ULA implementation with independent draws from
Eq.~\MainEqRef{eq:gmm-target} on the Kunitz memory. For each multiplicity ratio,
1{,}640 exact and 1{,}640 ULA states were decoded through the same PCA inverse
map and residue-wise argmax. The analytic designated-component probability
$f_{\mathrm{eff}}$ is reported separately from MAP-designated basin occupancy,
a hard assignment based on the maximum-posterior component
(Table~\ref{tab:gmm-benchmark}). The latter need not equal the mixture weight when
Gaussian components overlap.

\begin{table}[ht]
  \centering
  \caption{Exact Gaussian-mixture versus finite-step ULA sampling on Kunitz.
  Exact and ULA columns use equal sample counts and the same $\beta^{*}$.
  ``MAP des.'' is the fraction assigned to a designated maximum-posterior
  component and is not the latent mixture weight $f_{\mathrm{eff}}$.
  AA KL is the amino acid composition divergence between the decoded exact
  and ULA libraries.}
  \label{tab:gmm-benchmark}
  \small
  \resizebox{\textwidth}{!}{%
  \begin{tabular}{cccccccc}
    \toprule
    $\rho$ & $\beta^{*}$ & $f_{\mathrm{eff}}$
      & MAP des., exact & MAP des., ULA
      & P1 K/R, exact & P1 K/R, ULA & AA KL \\
    \midrule
    1   & 4.27 & 0.323 & 0.307 & 0.281 & 0.398 & 0.385 & $1.40\times10^{-4}$ \\
    10  & 4.93 & 0.827 & 0.894 & 0.887 & 0.515 & 0.499 & $1.25\times10^{-4}$ \\
    500 & 8.78 & 0.996 & 0.997 & 0.999 & 0.604 & 0.597 & $1.19\times10^{-4}$ \\
    \bottomrule
  \end{tabular}
  }
\end{table}
